\documentclass[11pt]{article}
\usepackage[margin=1in]{geometry}
\usepackage[T1]{fontenc}
\usepackage[utf8]{inputenc}
\usepackage{lmodern}
\usepackage{amsmath}
\usepackage{enumitem}
\usepackage{amsfonts}
\usepackage{amssymb}

\setlist[itemize]{noitemsep, topsep=2pt}
\setlist[enumerate]{noitemsep, topsep=2pt}

\title{TickRange Accrual Notes Specification}
\author{}
\date{\today}

\begin{document}
\maketitle
\tableofcontents

 
RANGE-ACCRUAL-NOTE::
	underlying: feeGrowthInside(liquidityPosition);
	strikeInterval: TickRange;
	coupon(period):  inStrikeInterval(price(swap)) ? fee*swap : 0x00;  
	observationPeriods: <NUMBER_OF_BLOCKS>, <...>
	accrualFrequency: <2 BLOCKS>, <N BLOCKS>, ..., <TIME_FREQUENCY>
	payoff(RANGE_ACCRUAL_NOTES) -> feeGrowthInside


        RANGE-ACCRUAL-NOTE::BOND

        BOND::
	   principal: 

 \subsection{Requirements}

R1: Accrues cupon in each period (Already fullfilled by Uniswap design)
R2: Pays out the coupon at the end of the block (Start of the new block)
R3: strikeInterval can be changed
R4: The accrued cupon can change across periods (Dynamic Fees)
R5: entitles the holder to coupon payments
R6: entitles the holder the payment of the principal
R7: calculates the payoff requireing the whole trajectory of the underlying price
 \subsection{Examples}
 \subsubsection{ 6 block wsETTH/WBTC feeGrowth}
....

 \subsection{Single-Period}

 label the instrument \texttt{TAN}.

 Let \(\Theta = \int_{\psi} \phi (\Delta (\psi)) d\psi\) be the representative LP position (\text{feeGrowthInside})

 We define a observationPeriods \(N\) can be modeled as number of blocks, number of transactions, or even time related (but discrete)

 Compared to Range Accrual notes \(\Theta\) will be the equivalent to the \textbf{coupons}.

 Given a sequence of \(n\) swaps \(\Delta_i\)_{i=1}^n that happened within the observationPeriod \(N\)
 which realized exhange rate induced sequence is \(P (\Delta_i)\) we define the \textbf{singe-period} tick range accrual payoff as:

 \begin{align*}
   \Pi^{\texttt{TAN}} = \Theta \cdot \frac{\sum^n \mathbb{1}_{ \psi_b \leq P(\Delta_i) \leq \psi_a}}{N}
 \end{align*}


 
 \subsection{Multi-Period}

 Given \(N =(N_j)^d\) then:

 \begin{align*}
     \Pi^{\texttt{TAN}} = \sum^j \Theta_j \cdot \frac{\sum_j \sum_i \mathbb{1}_{ \psi_b \leq P(\Delta_{ij}) \leq \psi_a}}{N_j} 
  \end{align*}


 
 \section{Extensions}

Extensions(RANGE_ACCRUAL_NOTE)::
	Accrual_Decrual(RANGE_ACCRUAL_NOTE) -> RANGE_ACCRUAL_NOTE
		updateAccrualFrequency(RANGE_ACCRUAL_NOTE, Frequency)
	rule: price(Accrual_Decrual(RANGE_ACCRUAL_NOTE) < price(RANGE_ACCRUAL_NOTE);
	
	TargetRedemptionNote(RANGE_ACCRUAL_NOTE) -> RANGE_ACCRUAL_NOTE
		setRedemptionTargetCap(RANGE_ACCRUAL_NOTE, Cap)

	Barriers(RANGE_ACCRUAL_NOTE) -> RANGE_ACCRUAL_NOTE
		setKnockInBarrier(RANGE_ACCRUAL_NOTE, Barrier)
		setKnockOutBarrier(RANGE_ACCRUAL_NOTE, Barrier)

	rule: price(Barriers(RANGE_ACCRUAL_NOTE)) < price(RANGE_ACCRUAL_NOTE)

	Callable(RANGE_ACCRUAL_NOTE) -> Option(RANGE_ACCRUAL_NOTE)
		mint(
			OptionType(Option(RANGE_ACCRUAL_NOTE))= CALL,
			OptionPosition(Option(RANGE_ACCRUAL_NOTE)) = SHORT
			strikePrice(Option(RANGE_ACCRUAL_NOTE))
		)
	 rule: price(Callable(RANGE_ACCRUAL_NOTE)) < price(RANGE_ACCRUAL_NOTE)
	 rule: Automatically repaid to the holder if certain conditions are met

	 BasketUnderlier(RANGE_ACCRUAL_NOTE) -> RANGE_ACCRUAL_NOTE
	 	updateUnderlying(RANGE_ACCRUAL_NOTE, MultiAssetVault (IERC7575))

	FloatingRangeAccrual(RANGE_ACCRUAL_NOTE) -> RANGE_ACCRUAL_NOTE
		updateCouponRule(RANGE_ACCRUAL_NOTE, TradingFeePolicy)



\end{document}
